% !TEX root = ../main.tex
\documentclass[../main.tex]{subfiles}

\begin{document}
\title{
  Galois Groups of Trinomials with Small Coefficients
}
\maketitle
% \begin{abstract}
%   We show that the Galois group over $\Q$ of any factor of $x^n + ax + b$ is the full symmetric group for all $n \in \N$ and for all pairs $\left( a, b \right)$ in the list \InlineProvenCases
  
% \end{abstract}
\localtableofcontents
\section{Introduction}
By Hilbert's irreducibility theorem \cite{Hilbert}, we expect a general rational polynomials of degree $n$ to have Galois group $S_n$ over $\Q$.
Nevertheless, giving explicit examples of such polynomials for all $n$ remained a difficult problem.
Only 38 years after Hilbert's work was published, Schur \cite{schur1931affectless} showed that the Laguerre polynomials 
\[
  L_n^{\left( 0 \right)}\left( -x \right) = \sum_{k=0}^n \binom{n}{k} \frac{x^k}{k!}
\]
have Galois group $S_n$ for every $n$.
Later, the same was proven for other families of generalized Laguerre polynomials, such as the Hermite polynomials, $L^{\left( \pm 1/2 \right)}_n$ \cite{schur1931affectless}, and the reverse Bessel polynomials, $L^{\left( -2n-1 \right)}_n$ \cite{grosswald1978bessel} \cite{filaseta2002irreducibility}.

Note that in all of these examples, the height of the coefficients tend to infinity rather fast.
Recent results show that the probability that the Galois group of a polynomial is either $A_n$ or $S_n$, tends to 1 as $n$ tends to infinity even if the coefficients are distributed according to certain finitely supported probability measures.
For example Bary-Soroker, Koukoulopoulos and Kozma \cite{Bary_Soroker_2023} proved that for a family of probability measures that includes the uniform distribution on any set of at least $35$ consecutive integers.
Assumeing some version of the Great Riemann Hypothesis, Breuillard and Varj\'{u} \cite{Breuillard_2019} proved that the same for another family of probability measures that includes the uniform distribution on ${0, 1}$, thus conditionally resolving a conjecture of Odlyzko and Poonen.

However, it appears that the only known example with a bounded height is the family of Selmer trinomials $f_n\left( x \right) = x^n - x - 1$ , which was shown to be irreducible by Selmer \cite{SelmerIrreducibility1956} and to have Galois group $S_n$ by Nart and Vila \cite{NartVila1979}. 
To achieve that, they first showed that all the inertia groups of $f_n$ are either trivial or generated by a transposition.
They then utilized Hermite's Theorem to prove that the Galois group of $f_n$ over $\Q$ is generated by its inertia groups, and deduced that it is $S_n$ since any transitive permutation group generated by transpositions is $S_n$.

It is natural to try to apply a similar strategy for $f_n\left( x \right) = x^n + ax + b$, or its irreducible factors, for some other values of $a,b \in \Z$.
Nart and Vila \cite{NartVila1979} themselves and later Osada \cite{Osada1987TheGG} indeed generalized the above method to obtain some sufficient criteria on $a,b$ and $n$ for the Galois group of $f_n$ to be $S_n$, assuming its irreducibility.
Yet, each of these criteria includes some congruence condition on $n$, leaving out a positive density of natural numbers as exceptions.
Maybe the most instructive example for how these congruence conditions arise is the  reducible famiy $f_n = x^n - 2x + 1$, which prompted our own interest in the subject due to a recent result of Bary-Soroker, Shusterman, and Zannier \cite[Appendix A]{Soroker_appendix_2017}.
In this case, $f_n$ factors as the product of $x-1$ and the polynomial $g\left( x \right) = x^{n-1} + x^{n-2} + \cdots + x - 1$.
The latter polynomial was shown to be irreducible by Miles \cite{Miles1960} and Wolfram \cite{Wolfram1998}, which were in fact interested in its close relative, the generalized Fibonacci polynomial $x^{n-1} - x^{n-2} - \cdots - x - 1$.
Martin \cite{MARTIN2004213} succeded, using methods similar to those used in the classical case of $x^n - x - 1$, to show that the Galois group of $g$ is $S_{n-1}$ for any odd value of $n$.
The reason his proof fails for even values of $n$, is that in this case 
\[
  x^n - 2x + 1 \equiv x^n - 1 \equiv \left( x^{n/2} - 1\right)^2 \mod 2
\]
and thus the inertia group of $f_n$ over $\Q_2$ is no longer generated by a transposition.

In this paper we present a proof that bypasses this problem for the mentioned case and additional ones.
\editorial{Here I will present the results, which I skip for now until we will be certain what exactly they are. Maybe the rest of the paragraph will also be rewritten, to highlight some lemmas that might be useful for other people.}
The previous arguments, both for the cases of generalized Laguerre polynomials, and for the case of $x^n - x - 1$, focused on local methods, which are insufficient for trinomials with bigger coefficients. 
Our new approach utilizes global bounds on the discriminant of number fields, like the Minkowski's bound and the stronger Odlyzko's bound \cite{OdlyzkoBoundsII}, to show that the "unproblematic" inertia groups, those that are indeed generated by transpositions, is already trnasitive, rendering the problematic ones dispensable for the proof.

% We will present the startegy of the proof that the Galois group is $S_n$ in the cases where $f$ is irreducible.
% The startegy for the reducible cases is similar.
% We will first show the key fact, for any prime number $p$ that does not divide $ab$, and an embedding of $\overline{\Q}$ into $\overline{\Q_p}$ the image of the absolute inertia group of $\Q_p$ in the Galois group of $f$ is either trivial or generated by a transposition.
% Assume towards contradiction that the Galois group of $f$ is not $S_n$.
% It is well known that any primitive permutation group that contains a transposition is symmetric.
% Therefore, either the Galois group of $f$ contains no transposition, in which case it follows from the key fact above that $f$ ramifies only above primes that divide $ab$, or the Galois group of $f$ is imprimitive.
% In the latter case, we will show that it follows that $\Q\left[ x \right] / \left( f \right)$ has a nontrivial subextension that is ramified only above primes that divide $ab$.
% Thus, in either case we will deduce that there is some subextension $\Q \subsetneq K \subseteq \Q\left[ x \right] / \left( f \right)$ that is ramified only above primes that divide $ab$.
% Finally, in each of the cases we will find a constant bound for the $p$-adic valuation of the root discriminant of $K$.
% In cases where the bound is good enough, the Odlyzko bound tells us that there are only finitely many number fields with root discriminant below the bound.
% In cases where the bound is even better, that list is enumerated in the LMFDB \cite{lmfdb}.
% If we are even luckier, there are very few number fields in this list that are ramified only above the primes that divide $ab$, and we can rule out by hand the possibility that each of these fields is contained in $\Q\left[ x \right] / \left( f \right)$.


% \subsection{Notations}
% Denote by $G$ the absolute Galois group $\Gal\left( \overline{\Q}/\Q \right)$, and for each prime $p$, denote by $I\left( p \right) \subseteq G$ the absolute inertia group at $p$, i.e., the embedding into $G$ of the inertia group $\Gal\left( \overline{\Q_p}/\Q_p^{\un} \right)$, where $\Q_p^{\un}$ is the maximal unramified extension of $\Q_p$ inside $\overline{\Q_p}$.

% For each polynomial $h$ and number field $K$, we will denote by $\d\left( h \right)$ and $\d\left( K \right)$ the discriminant of $h$ and of $K$ respectively and by $\rd\left( K \right)$ the root discriminant of $K$, i.e. $\rd\left( K \right) = \d\left( K \right)^{1/\left[ K : \Q \right]}$.

\section{Sufficient Condition }
\subsection{Setup and notations \label{subsection:strategy_setup}}
Let $a, b \in \Z - \left\{ 0 \right\}$ and $2 \leq n \in \Z$ and let $f\left( x \right) = x^n + ax + b$.
Let $g$ be an irreducible factor of $f$, let $L = \Q\left[ x \right] / \left( g \right)$, and denote by $\Lroot$ the image of $x$ in $L$.
$a,b,n,f,g,L$ will be fixed throughout this section and the next one.

For any field $\F$ fix a separable closure $\overline{\F}$ and denote by $G_\F$ its absolute Galois group.

Let $\ell$ be a prime not dividing $ab$.
Note that each embedding $\iota: \overline{\Q} \hookrightarrow \overline{\Q_\ell}$ induces an embedding $G_{\Q_\ell} \hookrightarrow G_\Q$.
Denote by $I\left( \iota \right)$ the image of the inertia group of $\Q_p$ under this embdding.
Let $H$ be the product of all the inertia groups $I\left( \iota \right)$ for all of the embeddings $\iota: \overline{\Q} \hookrightarrow \overline{\Q_\ell}$ and all of the primes $\ell$ not dividing $ab$.

\subsection{General Properties of $f$}
We will prove two lemmas that will be useful in a few places in the sequel. 
\begin{lemma}
  \label{lem:derivative_formula}
  If $\alpha$ is a root of $f$ in some field in which $b \neq 0$, then $\alpha$ is invertible and
  \[
    f'\left( \alpha  \right) = \left( 1 - n \right)a - nb\alpha^{-1}
  \]
\end{lemma}
\begin{proof}
  \begin{align*}
    f'\left( \alpha \right) = &
    n\alpha^{n-1} + a =
    n \frac{\alpha^n}{\alpha} + a = 
    n \frac{-a\alpha - b}{\alpha} + a =
     \left( 1 - n \right)a - nb\alpha^{-1} 
  \end{align*}
\end{proof}

\begin{lemma}
  \label{lem:multiple_root}
  If $\F$ is a field in which $ab \neq 0$, then $f$ is either separable over $\F$ or has $n-2$ simple roots and one double root in $\overline{\F}$.
\end{lemma}
\begin{proof}
  Let $\alpha \in \overline{\F}$ be a multiple root of $f$.
  Then, $\alpha$ is also a root of $f'$ and hence, since $b \neq 0$, it follows from Lemma \ref{lem:derivative_formula} that $\alpha \neq 0$ and that
  \[
    0 = f'\left( \alpha \right) = \left( 1 - n \right)a - nb\alpha^{-1}
  \]
  However, since $a$ is also non zero in $\F$, and since $n$ and $n-1$ cannot both be zero in $\F$, it follows that neither $n$ nor $n-1$ is zero in $\F$.
  Therefore, the equation above is linear with non zero coefficients and hence has a unique solution in $\overline{\F}$.
  Thus, this solution is the unique double root of $f$ in $\overline{\F}$ and it is only left to show that it is not a triple root of $f$. 
  However, the only root of $f''$ is zero which is not a root of $f$ and hence $f$ has no triple roots.  
\end{proof}

\subsection{Proof of the condition}

When we say that a group acts on a set $X$, as a single transposition, we mean that its image in $\Sym\left( X \right)$.
This terminology allows us to translate Lemma \ref{lem:multiple_root} into the following corollary which will be usefl in the proof of our condition:
\begin{corollary}
  \label{cor:inertia_group_acts_as_transposition}
  For any prime $\ell \not | ab$ and an embedding $\iota: \overline{\Q} \hookrightarrow \overline{\Q_\ell}$, the inertia group $I\left( \iota \right)$ acts on $\Hom_{\Q}\left( L, \overline{\Q} \right)$ either trivially or as a single transposition.
\end{corollary}
\begin{proof}
  It suffices to show the lemma for the $\Q$-algebra $\Q\left[ x \right] / \left( f \right)$ instead of $L$.
  Since $f$ is integral, it follows from \ref{lem:multiple_root} applied to $\F = \F_\ell$.
\end{proof}
\begin{remark}
  During the above proof we showed that $\ell \mid a b \left( n - 1 \right)$, which is the first sentence in the proof of \cite[lemma 3]{Osada1987TheGG}, and hence could have used Osada's lemma to conclude our own. 
\end{remark}

The second lemma that our proof relies on is the following global one:
\begin{lemma}
  \label{lem:transitive_action}
  If $H$ acts non transitively on $\Hom_{\Q}\left( L, \overline{\Q} \right)$, then there is some subextension $\Q \subsetneq K \subseteq L$ that is ramified only above primes that divide $ab$.
\end{lemma}
\begin{proof}
  We will show it using the famous fact that the functor $\Hom\left( -, \overline{\Q} \right)$ is a contravariant equivalence between the category of \'{e}tale $\Q$-algebras and the category of finite $G_\Q$-sets.
  Since $H$ is normal in $G_\Q$, $G_\Q$ acts on $\Hom_{\Q}\left( L, \overline{\Q} \right) / H$  through $G_\Q/H$.
  In particular, the quotient map $\Hom_{\Q}\left( L, \overline{\Q} \right) \to \Hom_{\Q}\left( L, \overline{\Q} \right) / H$ is an epimorphism in the category of $G_\Q$-sets.
  Therefore, there is a corresponding embedding of fields $K \hookrightarrow L$.
  
  Since $H$ acts non transitively on $\Hom_{\Q}\left( L, \overline{\Q} \right)$, $\Hom_{\Q}\left( L, \overline{\Q} \right) / H$ has at least two elements and hence $K \neq \Q$.
  Additionally, since $\Hom_{\Q}\left( L, \overline{\Q} \right) / H$ is fixed by $H$, it follows that $K$ is fixed by $H$ and is hence ramified only above primes that divide $ab$.
\end{proof}

The previous two lemmas together give us the following corollary:
\begin{thm}
  \label{thm:strategy}
  If there is no subextension $\Q \subsetneq K \subseteq L$ that is ramified only above primes that divide $ab$, then the Galois group of $g$ over $\Q$ is $S_{\deg\left( g \right)}$.
\end{thm}
\begin{proof}
  By \cite[Lemma 5]{Osada1987TheGG}, any transitive permutation group generated by transpositions is symmetric.
  Since $H$ is generated by subgroups that act on $\Hom_{\Q}\left( L, \overline{\Q} \right)$ either trivially or as a transposition according to Lemma \ref{cor:inertia_group_acts_as_transposition}, it follows that the image of $H$ in $\Sym\left( \Hom_{\Q}\left( L, \overline{\Q} \right) \right)$ is generated by trnaspositions.
  However, this image is also transitive as a result of Lemma \ref{lem:transitive_action} and hence it is the full symmetric group.
\end{proof}
We shall then prove that for some values of $a,b$, the discriminant of $L$ is too small for such subextension $K$ to exist.

\section{Effective Sufficient Condition}
The main aim of this section is proving Theorem \ref{thm:criterion} \editorial{which will have appeared in the Introduction}.
We will do it by proving some bounds on the $p$-adic valuation of the discriminant of $L$. Some of these bounds will be useful in their own right in the next section.

\subsection{Setup and notations}
Let $a,b,n,f,g,L,\Lroot$ be as defined in subsection \ref{subsection:strategy_setup}.
Let $\d\left( L \right)$ be the discriminant of $L$ and $\rd\left( L \right)$ be the root discriminant of $L$, i.e. $\rd\left( L \right) = \d\left( L \right)^{1/\left[ L : \Q \right]}$.
Let $p$ be a prime dividing $ab$.
$p$ will be fixed throughout this subsection and subsection \ref{subsection:discriminant_bounds}.
Let $\pi = \pi_p: \O_{\overline{\Q_p}} \to \overline{\F_p}$ be the cannonical quotient map from the integral closure of $\Z_p$ inside $\overline{\Q_p}$.
Let $v_p: \overline{\Q_p} \to \Q$ be the unique extension of the $p$ adic valuation on $\Q_p$.
Fix an embedding of $\overline{\Q}$ into $\overline{\Q_p}$.
This embedding induces a valuation on $\overline{\Q}$ which will be denoted by $v_p$ as well. Use the standard notation for the $p$-adic norm $\pnorm{\cdot} = p^{-v_p\left( \cdot \right)}$.
For each embedding $\Lembedding: L \to \overline{\Q}$, we will denote by $\Lembedding\left( L \right)\Q_p$ the compositum of $\Lembedding\left( L \right)$ and $\Q_p$ inside $\overline{\Q_p}$, and by $\D\left( \Lembedding \right)$ the different of $\Lembedding\left( L \right)\Q_p$ over $\Q_p$.
We press that in contrast to the previous section, where the embeddings $\iota: \overline{\Q} \hookrightarrow \overline{\Q_\ell}$ vary, in this section the embedding $\overline{\Q} \hookrightarrow \overline{\Q_p}$ is fixed and the embeddings $\Lembedding: L \to \overline{\Q}$ vary.

\subsection{$p$-adic bounds on the discriminant \label{subsection:discriminant_bounds}}
To bound the valuation of the discriminant of $L$, we first bound the valuation of $\Deta$ for each embedding $\Lembedding: L \to \overline{\Q_p}$.
To bound this valuation in turn, we wish we could bound the valuation of $f'\left( \Lembedding\left( \Lroot \right) \right)$ for each $\Lembedding$.
Instead we bound the valuation of $f'\left( \Lembedding\left( \Lroot \right) \right)$ for most of the embeddings $\Lembedding$, and characterize the remaining ones.
\begin{lemma}
  \label{lem:valuation_of_derivative}
  Assume that $n > v_p\left( b \right)$.
  For every homomorphism $\Lembedding: L \to \overline{\Q_p}$ one of the following holds: 
  \begin{enumerate}
    \item \label{case:good} $v_p\left( f'\left( \Lembedding\left( \Lroot \right) \right) \right) \leq \maxpval$.
    \item \label{case:single_slope} $v_p\left( b \right) = 0$, $v_p\left( a \right) = v_p\left( n \right)$ and $\Lembedding$ is one of the $\ppart{a}$ embeddings for which $\pi\left( \Lembedding\left( \Lroot \right) \right) = \pi\left( \frac{nb}{\left( 1 - n \right)a} \right)$.
    \item \label{case:longer_slope} $v_p\left( a \right) = 0$, $v_p\left( b \right) = v_p\left( n - 1 \right)$ and $\Lembedding$ is one of the $\ppart{b}$ embeddings for which $\pi\left( \Lembedding\left( \Lroot \right) \right) = \pi\left( \frac{nb}{\left( 1 - n \right)a} \right)$.
    \item \label{case:shorter_slope} $v_p\left( b \right) > v_p\left( a \right)$, $v_p\left( n-1 \right) = 0$ and $\Lembedding$ is the unique embedding for which $v_p\left( \Lembedding\left( \Lroot \right) \right) = v_p\left( b \right) - v_p\left( a \right)$.
  \end{enumerate}
\end{lemma}
\begin{proof}
  Assume first that 
  \[
    v_p\left( n - 1 \right) + v_p\left( a \right) \neq 
    v_p\left( n \right) + v_p\left( b \right) - v_p\left( \Lembedding\left( \Lroot \right) \right).
  \]
  Then by Lemma \ref{lem:derivative_formula} 
  \begin{align*}
    v_p\left( f'\left( \Lembedding\left( \Lroot \right) \right) \right) = 
    & v_p\left( \left( 1 - n \right)a - nb\Lembedding\left( \Lroot \right)^{-1} \right) = \\ 
    & \min{\left\{ 
      v_p\left( n - 1 \right) + v_p\left( a \right), 
      v_p\left( n \right) + v_p\left( b \right) - v_p\left( \Lembedding\left( \Lroot \right) \right) 
    \right\}} \leq \\
    & \min{\left\{ 
      v_p\left( n - 1 \right) + v_p\left( a \right), 
      v_p\left( n \right) + v_p\left( b \right) 
    \right\}}.
  \end{align*}
  However, since either $v_p\left( n \right) = 0$ or $v_p\left( n - 1\right) = 0$, it then follows that indeed 
  \[
    v_p\left( f'\left( \Lembedding\left( \Lroot \right) \right) \right) \leq
    \max\left\{ v_p\left( a \right), v_p\left( b \right) \right\} = 
    \maxpval
  \]
  and hence item $\ref{case:good}$ holds.

  Therefore, we can assume that
  \begin{equation}
    \label{eq:valuation_inequality}
    v_p\left( f'\left( \Lembedding\left( \Lroot \right) \right) \right) > \maxpval
  \end{equation} 
  and as a result that
  \begin{equation}
    \label{eq:valuation_equality}
    v_p\left( n - 1 \right) + v_p\left( a \right) = 
    v_p\left( n \right) + v_p\left( b \right) - v_p\left( \Lembedding\left( \Lroot \right) \right)
  \end{equation}
  In particular, it follows that $v_p\left( \Lembedding\left( \Lroot \right) \right) \in \Z$.
  Therefore, $\Lembedding\left( \Lroot \right)$ must correspond to a segment of the Newton polygon of $f$ with respect to $v_p$ whose slope is an integer.
  Recalling that $n > v_p\left( b \right)$, one can see that, as illustrated in figure \editorial{add a figure and a reference to it}, it can possibly happen in one of only three ways:
  \begin{itemize}
    \item $v_p\left( b \right) = 0 = v_p\left( \Lembedding\left( \Lroot \right) \right)$. 
    We will show that item $\ref{case:single_slope}$ holds.
    Indeed, Equation \ref{eq:valuation_equality} specializes to 
    \[
      v_p\left( n - 1 \right) + v_p\left( a \right) = v_p\left( n \right) + v_p\left( b \right)
    \]
    The left hand side is positive since $p$ must divide $a$, and hence the right hand side is positive, and hence $p$ divides $n$.
    Therefore, $v_p\left( n - 1 \right) = 0$ and hence the only non zero terms in Equation \ref{eq:valuation_equality} are $v_p\left( n \right)$ and $v_p\left( a \right)$ and hence $v_p\left( n \right) = v_p\left( a \right)$.
    By inequality \ref{eq:valuation_inequality} and Lemma \ref{lem:derivative_formula}, 
    \[
      v_p\left( a \right) < 
      v_p\left( f'\left( \Lembedding\left( \Lroot \right) \right) \right) = 
      v_p\left( a \right) + v_p\left(  
        \left( 1 - n \right)a\pnorm{a} - n\pnorm{n}b\Lembedding\left( \Lroot \right)^{-1}
      \right)
    \]
    and hence 
    \[
      \pi\left( \Lembedding\left( \Lroot \right) \right) = 
      \pi\left( \frac{nb}{\left( 1 - n \right)a} \right)
    \]
    Since 
    \[
      x^n + ax + b \equiv
      \left( x^{\left( \ppart{a} \right)} \right)^{n\pnorm{a}} + b \mod p
    \]
    the fiber of each root of $f$ in $\overline{\F_p}$ contains exactly $\ppart{a}$ among the roots of $f$ in $\overline{\Q_p}$. 
    Consequently, there are exactly $\ppart{a}$ homomorphisms $\Lembedding: L \to \overline{\Q_p}$ for which $\pi\left( \Lembedding\left( \Lroot \right) \right) = \pi\left( \frac{nb}{\left( 1 - n \right)a} \right)$.
    
    \item $v_p\left( b \right) > v_p\left( a \right) = 0 = v_p\left( \Lembedding\left( \Lroot \right) \right)$.
    The analysis of this case is similar to the previous one, and we will show that item $\ref{case:longer_slope}$ holds.

    \item $v_p\left( \Lembedding\left( \Lroot \right) \right) = v_p\left( b \right) - v_p\left( a \right) > 0$.
    We will show that item $\ref{case:shorter_slope}$ holds.
    Substituting into Equation \ref{eq:valuation_equality} we get
    \[
      v_p\left( n - 1 \right) + v_p\left( a \right) = 
      v_p\left( n \right) + v_p\left( b \right) - \left( v_p\left( b \right) - v_p\left( a \right) \right) = v_p\left( n \right) + v_p\left( a \right)
    \]
    from which it follows that $v_p\left( n - 1 \right) = v_p\left( n \right)$ and hence $0 = v_p\left( n \right) = v_p\left( n - 1 \right)$.
  \end{itemize}
\end{proof}

% For each extension $K/F$ of field of fractions of Dedekind domains $B/A$, let $\D\left( K/F \right)$ denote the different of $B$ over $A$.
% For each rational integer $m$, denote $ = p^{v_p\left( m \right)}$.
As promised, we now utilize the previous lemma to bound the $v_p\left( \D_\eta \right)$ for each embedding $\Lembedding: L \hookrightarrow \overline{\Q_p}$.
\begin{lemma}
  \label{lem:different_valuation}
  Assume that $n > v_p\left( b \right)$.
  Then for each $\Lembedding: L \to \overline{\Q_p}$, one of the following holds:
  \begin{enumerate}
    \item $v_p\left( \Deta \right) \leq \maxpval$.
    \item $v_p\left( a \right) + v_p\left( n - 1 \right) = v_p\left( b \right) + v_p\left( n \right)$, $\Lembedding$ is one of the only $\maxppart$ embeddings for which $\pi\left( \Lembedding\left( \Lroot \right) \right) = \pi\left( \frac{nb}{\left( 1 - n \right)a} \right)$, and 
    \begin{align*}
      v_p\left( \Deta \right) \leq & 
      \diffbnd 
    \end{align*}
  \end{enumerate}
\end{lemma}
\begin{proof}
  We use the same division into cases as in the previous lemma.
  In the first case, 
  \begin{equation*}
    v_p\left( \Deta \right) \leq
    v_p\left( f'\left( \Lembedding\left( \Lroot \right) \right) \right) \leq 
    \maxpval
  \end{equation*}
  and hence the first item of the corollary holds.

  Assume the third case holds, i.e., $v_p\left( b \right) > v_p\left( a \right)$, $v_p\left( n-1 \right) = 0$ and $\Lembedding$ is the unique embedding for which $v_p\left( \Lembedding\left( \Lroot \right) \right) = v_p\left( b \right) - v_p\left( a \right)$.
  We will show that in this case the first item of the corollary also holds.
  Indeed, since $\Lembedding$ is the only embedding for which $v_p\left( \Lembedding\left( \Lroot \right) \right)$, it follows in particular that $\Lembedding\left( \Lroot \right)$ is not conjugate over $\Q_p$ to any of its conjugates over $\Q$, and hence belongs to $\Q_p$ and hence $v_p\left( \Deta \right) = 0$.

  Assume that either the second or the fourth case holds.
  We will show that in both cases the second item of the corollary holds.
  Indeed, in both cases 
  \[
    v_p\left( a \right) + v_p\left( n - 1 \right) = v_p\left( b \right) + v_p\left( n \right)
  \]
  and $\Lembedding$ is one of the exactly $\maxppart$ embeddings for which $\pi\left( \Lembedding\left( \Lroot \right) \right) = \pi\left( \frac{nb}{\left( 1 - n \right)a} \right)$.
  Denote by $e\left( \Lembedding \right)$ the ramification index of $\Lembedding\left( L \right)\Q_p$ over $\Q_p$.
  Since the ramification index of $\Lembedding\left( L \right)\Q_p$ over $\Q_p$ is the cardinality of $\Lembedding\left( \Lroot \right)$ under the action of the absolute inertia group of $\Q_p$, it follows that $$e\left( \Lembedding \right) \leq \maxppart$$ and in particular that 
  \[
    v_p\left( e\left( \Lembedding \right) \right) \leq \maxpval
  \]
  We can combine these two inequalities with the well known bound (e.g. \cite[p. 58]{SerreLocalFields}) 
  \[
    v_p\left( \Deta \right) \leq 
    1 - e\left( \Lembedding \right)^{-1} + v_p\left( e\left( \Lembedding \right) \right)
  \]
  to deduce that
  \begin{align*}
    & v_p\left( \Deta \right) \leq
     \diffbnd
  \end{align*}
\end{proof}

We now utilize this lemma to finally obtain a bound on the $p$-adic valuation of the discriminant of $L$:
\begin{corollary}
  \label{cor:discriminant_bound}
  Assume that $n > v_p\left( b \right)$.
  Then
  \[
    v_p\left( \d\left( L \right) \right) \leq 
    \discbnd
  \]
\end{corollary}
\begin{proof}
  By the last lemma,\editorial{Check this again!}  
  \begin{align*}
    & v_p\left( \d\left( L \right) \right) = \\
    &  \sum_{\eta: L \hookrightarrow \overline{Q}} v_p\left( \Deta \right) \leq \\
    & \left( \deg L - \maxppart \right)\maxpval + 
    \maxppart \left( \diffbnd \right) \leq \\
    & \discbnd 
    % & \deg L\maxpval + 2a^2b^2
  \end{align*}
\end{proof}

Before continuing with the proof of Theorem \ref{thm:criterion}, we will record a slight improvement to the previous corollary that will be useful in the analysis of some of the cases in the next section:
\begin{corollary}
  \label{cor:discriminant_bound_Eisenstein}
  If $f$ is Eisenstein at $p$ and $n > v_p\left( b \right)$, then
  \[
    v_p\left( \d\left( L \right) \right) \leq \deg L \cdot \maxpval
  \]
\end{corollary}
\begin{proof}
  Since the $p$-adic valuations of both $a$ and $b$ are positive, it follows that in Lemma \ref{lem:valuation_of_derivative} only the first and third case might hold, and hence in Lemma \ref{lem:different_valuation} only the first case might hold, and hence the sum in the proof of Corollary \ref{cor:discriminant_bound} is bounded by $\deg L \cdot \maxpval$.
\end{proof}

\subsection{Proof of the effective condition}

Finally we can prove a sufficient condition for the Galois group of $L$ to be the full symmetric group:
\begin{thm}
  \label{thm:criterion}
  Assume that there are constants $C_p > \maxpval$ for each prime $p|ab$, such that if a number field $K$ satisfies 
  \begin{enumerate}
    \item $v_p\left( \rd\left( K \right) \right) < C_p$ for each $p|ab$,
    \item $K$ is ramified only above primes dividing $ab$,
  \end{enumerate} 
  then $K$ is not contained in $L$.
  Assume that
  \[
    \deg L > \max_{p| ab}\left( \frac{\error}{C_p - \maxpval} \right)
  \] 
  Then the Galois group of $L$ over $\Q$ is $S_{\deg L}$.
\end{thm}
\begin{proof}
  We will show that the conditions of Corollary \ref{thm:strategy} hold.
  Let $K$ be a number field ramified only above primes dividing $ab$.
  Assume towards a contradiction that $K \subseteq L$.
  The inequality in the statement of the theorem is equivalent to the statement that 
  \[
    C_p > \maxpval + \frac{\error}{\deg L}
  \]
  for each $p|ab$.
  It follows, by the contrapositive of the first assumption in the theorem statement, that there is a prime $p|ab$ such that 
  \[
    C_p \leq v_p\left( \rd\left( K \right) \right) 
  \] 
  Therefore, if we denote by $d\left( L/K \right)$ the relative discriminant of $L$ over $K$, and by $N$ the norm from $K$ to $\Q$, then
  \begin{align*}
    & v_p\left( d\left( L \right) \right) = 
    v_p\left( \d\left( K \right)^{\left[ L:K \right]} N\left( \d\left( L/K \right) \right) \right) \geq \deg L \cdot v_p\left( \rd\left( K \right) \right) \geq \\
    & \deg L \cdot C_p > \deg L \cdot \left( \maxpval + \frac{\error}{\deg L} \right) = \\
    & \discbnd 
  \end{align*}
  which is a contradiction to corollary \ref{cor:discriminant_bound}.
\end{proof}

\section{Application Examples}
Carry over the notations and setup from the previous section.
We will now show how to apply Theorem \ref{thm:criterion} to prove that for some small values of $a,b$, the Galois group of $f$ is the largest that its factorization allows for all but at most finitely many values of $n$.
\subsection{$\lcm\left( a,b \right) = 1$}
If $\lcm\left( a,b \right) = 1$, then $f = x^n \pm x \pm 1$ and can always be reduced to either of the forms 
\[
  x^n - x - 1 \quad \text{or} \quad x^n + x + 1 
\]
by substituting $-x$ instead of $x$ if necessary.
Selmer \cite{SelmerIrreducibility1956} showed that the first form is always irreducible, while the second one is irreducible whenever $n \not \equiv 2 \mod 3$ and factors into $x^2 + x + 1$ times an irreducible factor otherwise.
When $x^n \pm x + 1$ is irreducible, the symmetry of its Galois group was established by \cite{NartVila1979} and \cite[Corollary 3]{Osada1987TheGG}.
It seems that this is the first time the reducible case is been analyzed as well:
\begin{corollary}
  \begin{align*}
    & \G\left( x^n  - x - 1 \right) = S_n \\
    & \G\left( x^n  + x + 1 \right) = 
    \begin{cases}
      S_n & n \not \equiv 2 \mod 3 \\
      S_{n-2} \times C_2 & n \equiv 2 \mod 3
    \end{cases}
  \end{align*}
\end{corollary}
\begin{proof}
  The Galois group of any divisor $g$ of $x^n \pm x \pm 1$ is $S_{\deg\left( g \right)}$ as a result of Theorem \ref{thm:strategy}:
  Indeed, there is no prime that divides $a b = \pm 1$ and the discriminant of any non-trivial extension of $\Q$ is divisible by some prime as a result of the Minkowski bound.
  It thus only remains to show that if $f = x^n + x + 1$ and $n \equiv 2 \mod 3$, then the Galois group of $x^2 + x + 1$ and the Galois group of the other irreducible factor of $f$, which shall be denoted by $g$, are disjoint.
  
  It suffices to show that that $\Q\left[ x \right] / \left( g \right)$ is not ramified at $3$, or equivalently, since $g$ is integral, that $g$ has only simple roots in $\F_3$.
  However, by Lemma $\ref{lem:multiple_root}$, $f$ has at least $n-2$ simple roots in $\overline{\F_3}$, and hence, since $x^2 + x + 1 \equiv \left( x - 1 \right)^2 \mod 3$, all of the $n-2$ roots of $g$ in $\overline{\F_3}$ are simple roots of $f$ and are in particular simple roots of $g$.
  \editorial{I could also show it using the explicit formula for $g$, but this argument felt nicer.}
\end{proof}

\subsection{Polynomial factorization}
We will record a few result about the factorization of $f$ for some values of $a,b,n$ that will be useful in the following examples.
\begin{lemma}
  \label{lem:irreducibility_b_1}
  Assume $\left| b \right| = 1$.
  \begin{enumerate}
    \item If $\left| a \right| > 2$ or $f = x^n + 2x - 1$, then $f$ is irreducible.
    \item If $f = x^n - 2x + 1$, then $f$ factors as $\left( x - 1 \right) g$ for some irreducible polynomial $g$.
  \end{enumerate}
  Up to substituting $-x$ instead of $x$, these are all the possible cases.
\end{lemma}
\begin{proof}
  In both cases, one can use Rouch\'{e}'s theorem to show that $f$ has a unique complex root inside the open unit disk.
  In the first case, $f$ has no roots on the unit circle, and hence is irreducible.
  Therefore, since its free term is $\pm 1$, any of its irreducible factors must have a root inside the unit circle, and hence $f$ can not have more than one irreducible factor.
  The second case is analogous to the case $x^n - 2x^{n-1} + 1$ which was proven by Martin \cite[Corollary 2]{MARTIN2004213}
\end{proof}

\begin{lemma}
  \label{lem:irreducibility_a_1}
  Assume that $\left| a \right| = 1$ and that $\left| b \right| = p^e$ for some positive integer $e$ and some prime $p$.
  \begin{enumerate}
    \item If $b > 2$ or $f = x^n - x + 2$, then $f$ is irreducible.
    \item If $f = x^n + x - 2$, then $f$ factors as $\left( x - 1 \right)$ times an irreducible polynomial.
  \end{enumerate}
  up to substituting $-x$ instead of $x$, these are all the possible cases.
\end{lemma}
\begin{proof}
  Again, by Rouch\'{e}'s theorem, $f$ has no roots inside the open unit disk, and in the first case, $f$ has no roots on the unit circle as well.
  In the second case one can show that the only root of $f$ on the unit circle is $1$.
  Let $f_0$ be $f$ in the first case and $f / \left( x - 1 \right)$ in the second case.
  In both cases, $f_0$ has no roots inside the closed unit disk, and we will show that $f_0$ is irreducible.

  Assume that $f_0$ factors as $g h$ for some non-constant polynomials $g,h \in \Z\left[ x \right]$.
  By analyzing the Newton polygon of $f$ with respect to $p$, it follows that $f$ has one root in $\overline{\Q_p}$ of valuation $e$ and $n-1$ roots of valuation $0$.
  Therefore, without loss of generality, we can assume that $g$ has only roots of valuation $0$.
  In particular, it follows that the free term of $g$ is $\pm 1$ and hence $g$ has a root in $\overline{\Q}$ of absolute value not greater than $1$ and that's a contradiction.
\end{proof}

\subsection{$\lcm\left( a,b \right) = 2$}
Assume that $\lcm\left( a,b \right) = 2$.
According to the LMFDB \cite{lmfdb}, the only number field with root discriminant not greater than $2$ that is ramified only above $2$ is $\Q\left( i \right)$.
The next real number realized as the root discriminant of a number field ramified only above $2$ is $\sqrt{8}$.
Therefore, once we rule out the possibility that $\Q\left( i \right)$ is contained in $\Q\left[ x \right] / \left( f \right)$ for any $n\geq 3$, we can apply Theorem \ref{thm:criterion} with $C_2 = \sqrt{8}$ which would hence apply whenever 
\[
  \deg L > \frac{\error}{\sqrt{8} - 2} = \frac{1}{\sqrt{8} - 2} < 1
\]
i.e., for all $n \geq 3$.
Therefore, we can conclude the following characterization of the Galois group of $f$ for all $n \geq 3$:
\begin{corollary}
  If $\lcm\left( a,b \right) = 2$ then up to substituting $-x$ instead of $x$, either of the following holds:
  \begin{enumerate}
    \item $f = x^n - 2x + 1, x^n + x - 2$ and $\G\left( f \right) = S_{n-1}$.
    \item $G\left( f \right) = S_n$ 
  \end{enumerate}
\end{corollary}
\begin{proof}
  For $n = 2$ the validity of the corollary can be checked manually.
  Assume that $n \geq 3$. 
  By Lemma \ref{lem:irreducibility_b_1} and Lemma \ref{lem:irreducibility_a_1}, in the first case $f$ factors as $\left( x - 1 \right)$ times an irreducible polynomial, and in the second case $f$ is irreducible. 
  Therefore, by the discussion above, we only need to show that $\Q\left( i \right)$ is not contained in $\Q\left[ x \right] / \left( f \right)$.
  Indeed, one of the following holds:
  \begin{itemize}
    \item $\left| a \right| = \left| b \right| = 2$. 
    Assume towards a contradiction that $\Q\left( i \right)$ is contained in $\Q\left[ x \right] / \left( f \right) = L$.
    Then the ramification index of $2$ in $L$ is $n$, and hence $\left( 1 + i \right)$ is ramified in $L$, and hence for each embedding $\Lembedding: L \to \overline{\Q_2}$, 
    \[
      v_2\left( \D_\iota \right) = 
      v_2\left( \D\left( \Q\left( i \right) / \Q \right) \right) + v_2\left( \D\left( \iota\left( L \right) \right) / \Q\left( i \right) \right) > v_2\left( \D\left( \Q\left( i \right) / \Q \right) \right) = 1 
    \]
    in contradiction to Corollary \ref{lem:different_valuation}.
    \item $f = x^n \pm 2x \pm 1$. In this case $f$ has a real root, and hence $\Q\left[ x \right] / \left( f \right)$ can not contain a totally imaginary number field as $\Q\left( i \right)$.
    \item $f = x^n \pm 2x \pm 1$. 
  \end{itemize}
\end{proof}
\end{document}